% Applied Software Development Lab -- Lab Handout 2: Blinking LED
% Copyright (c) Michael Lehn. Licensed under CC BY-SA 4.0.
%
% Region coordinates (bp from the top of the page) follow the page layout of
% the original German handout "Laborblatt 2".
\documentclass{aslab}
\sethandout{2}{Blinking LED}
\graphicspath{{../figures/handout02/}{../figures/shared/}}
\usepackage{pifont}
\usepackage{luacolor}
\usepackage{lua-ul}

\makeatletter
%% Lists used in this handout --------------------------------------------
%% "A)" lists with the text at +18bp
\newlist{alphatasks}{enumerate}{1}
\setlist[alphatasks]{label={\Alph*)},leftmargin=18bp,labelsep=*,align=left,
  labelindent=0pt}
%% plain "1." lists (regular weight)
\newlist{numlist}{enumerate}{1}
\setlist[numlist]{label={\arabic*.},leftmargin=18bp,labelsep=*,align=left,
  labelindent=0pt}
%% small bullet lists with the text at +9bp
\newlist{dotlist}{itemize}{1}
\setlist[dotlist]{label={\textbullet},leftmargin=9bp,labelsep=*,align=left,
  labelindent=0pt}
%% arrow lists with the text at +18bp
\newcommand*\aslabarrow{\ding{220}}
\newcommand*\aslabbigarrow{{\fontsize{13.2bp}{12bp}\selectfont\ding{220}}}
\newlist{arrowlist}{itemize}{1}
\setlist[arrowlist]{label={\aslabbigarrow},leftmargin=18bp,labelsep=*,
  align=left,labelindent=0pt}

%% Assembly listings with syntax colors ----------------------------------
%% Inside asmlisting, every space is a fixed 0.6em and every line break
%% ends a line, so the columns of the original listings are preserved.
\newcommand*\aN[1]{\textcolor{codeorange}{#1}}% line numbers, labels
\newcommand*\aI[1]{\textcolor{codegreen}{#1}}% instructions
\newcommand*\aK[1]{\textcolor{codered}{#1}}% constants, I/O addresses
\newcommand*\aR[1]{\textcolor{codecyan}{#1}}% registers, jump targets
\newcommand*\aC[1]{\textcolor{codepurple}{#1}}% comments
\newcommand*\asm@space{\leavevmode\hbox to 0.6em{}}
{\catcode`\ =\active\catcode`\^^M=\active%
\gdef\asm@setup{\catcode`\ =\active\def {\asm@space}%
\catcode`\^^M=\active\def^^M{\par}}}%
\newenvironment{asmlisting}{\par\ttfamily\asm@setup}{\par}

%% Absolute drawing in page coordinates (bp, y downwards) ------------------
\newcommand\pagedraw[1]{\placeat{0}{0}{\begin{tikzpicture}[x=1bp,y=-1bp]
  \useasboundingbox (0,0) rectangle (1,1);#1\end{tikzpicture}}}
\definecolor{tablegray}{HTML}{F5F5F5}
\definecolor{tableline}{HTML}{A5A5A5}
\definecolor{hlgray}{HTML}{E8E8E8}
\tikzset{tline/.style={draw=tableline,line width=0.5bp}}

%% Highlighted text (gray background as in the original), also across
%% line breaks: 10.3bp above and 4bp below the baseline
\newcommand\hltext[1]{\underLine[top=10.3bp,bottom=-4bp,color=hlgray]{#1}}
%% Table cells: 10bp text on a 12bp pitch, ragged right
\newcommand*\cell{\small\baselineskip=12bp\raggedright}
%% Keep the baseline pitch fixed even next to tall inline code boxes
\lineskiplimit=-\maxdimen
\makeatother

\begin{document}
%% ===================================================================== 1 ==
\maketitleblock

\begin{region}{124.0}
\heading{Learning objectives}
\end{region}

\begin{region}{149.0}
In the last lab handout, you created machine code by hand. Today we show
how the same machine code can be generated automatically from an assembly
program -- and then flashed to the microcontroller. You will also practice
working with the terminal.
\end{region}

\begin{region}{199.0}
\heading{Step 1: Working with the terminal}
\end{region}

\begin{region}[pitch=14.3]{227.5}
\textbf{Goal:} You will learn to find your way around the terminal, to
clone a repository with Git, and to navigate the file system.
\end{region}

\begin{region}[pitch=13.2]{269.4}
We start with two terminals side by side -- one for editing, one for
running commands.
\end{region}

\placegraphic{56.7}{286.4}{458.6}{terminals-git-clone.png}

\begin{region}{495.2}
\subheading{Left terminal (editing):}
\end{region}

\begin{region}[pitch=12.2]{509.4}
\begin{alphatasks}
\item \textbf{Show the directory contents}\\
  With the command \code{ls}, you can list the files and subdirectories in
  the current directory. Since it should be empty, nothing is displayed.
\end{alphatasks}
\end{region}

\begin{region}[pitch=12.2]{546.7}
\begin{alphatasks}[start=2]
\item \textbf{Clone the Git repository}\\
  \code{git clone https://gitlab.uni-ulm.de/mlehn/avr-bare-metal.git}
  downloads\\ the code base that we use for this lab handout.
\end{alphatasks}
\end{region}

\begin{region}[pitch=12.0]{584.1}
\begin{alphatasks}[start=3]
\item \textbf{Show the directory contents again}\\
  Enter \code{ls} again. The directory \code{avr-bare-metal} should now be
  displayed -- on the lab computers in blue, because it is a directory.
  Regular files appear in black.
\end{alphatasks}
\end{region}

\begin{region}[pitch=12.3]{633.4}
\begin{alphatasks}[start=4]
\item \textbf{Show the current directory}\\
  With \code{pwd} (= print working directory), you can see which directory
  you are currently in. This should be \code{/home/asp} -- your home
  directory. The directory asp is a subdirectory of \code{home}, which in
  turn is located directly in the root directory \code{/}.
\end{alphatasks}
\end{region}

\begin{region}[pitch=12.3]{683.3}
\begin{alphatasks}[start=5]
\item \textbf{Change into the subdirectory}\\
  With \code{cd avr-bare-metal}, you change into the new directory. Use
  \code{pwd} to check that you are now in \code{/home/asp/avr-bare-metal},
  and list its contents with \code{ls}. You should see a subdirectory
  \code{asm}.
\end{alphatasks}
\end{region}

\begin{region}[pitch=12.3]{733.2}
\begin{alphatasks}[start=6]
\item \textbf{One directory deeper and opening a file}\\
  Now change into the subdirectory with \code{cd asm} and check again with
  \code{pwd} and \code{ls} that you are now in
  \code{/home/asp/avr-bare-metal/asm}. You should see the file
  \code{led_pin2.s} -- the assembly program. Open it in the editor with
  \code{nano led_pin2.s}.
\end{alphatasks}
\end{region}

\nextpage
%% ===================================================================== 2 ==
\begin{region}{69.4}
\subheading{Right terminal (compiling \& flashing)}
\end{region}

\begin{region}[pitch=12.3]{82.4}
\begin{alphatasks}
\item \textbf{Change into the same directory}\\
  The challenge: in the right terminal, you have to change manually into
  the same directory as in the left terminal. You can do this in one step
  with \code{cd avr-bare-metal/asm}, going down two directories at once.
  Alternatively, you can first run \code{cd avr-bare-metal} and then
  \code{cd asm}.
\end{alphatasks}
\end{region}

\begin{region}[pitch=12.1]{144.4}
\begin{alphatasks}[start=2]
\item \textbf{Display the content of the file}\\
  With \code{cat led_pin2.s}, you can display the content of the file --
  this is assembly code, a human-readable form of machine code. With the
  right tools, it can be turned into a so-called hex file that can be
  transferred to the microcontroller.
\end{alphatasks}
\end{region}

\placegraphic{74.7}{195.4}{445.9}{terminals-nano-cat.png}

\begin{region}[pitch=12.15]{399.4}
\begin{alphatasks}[start=3]
\item \textbf{Understand the Makefile and compile}\\
  The so-called Makefile is also located in the current directory. It
  contains instructions on how the HEX file is generated automatically from
  the assembly program. Display it with \code{cat Makefile}. Then start the
  build process with \code{make}.
\end{alphatasks}
\end{region}

\placegraphic{74.7}{438.4}{384.5}{make-output.png}

\begin{region}[x=74.7,right=0,pitch=12.2]{511.5}
If everything works, a file \code{led_pin2.hex} is generated. It contains
the machine code that you created by hand in the last lab handout. Look at
its content with \code{cat led_pin2.hex}; it should look exactly like this:
\end{region}

\begin{region}[x=74.7,right=0,pitch=12.3]{560.2}
\code{:0C000000529A5A9A000000005A98FBCF58}\\
\code{:00000001FF}
\end{region}

\begin{region}[x=74.7,right=0]{596.5}
This is exactly the machine code from Lab Handout~1 that switches pin~2
between HIGH and LOW every 250\,ns.
\end{region}

\begin{region}[pitch=12.15]{621.5}
\begin{alphatasks}[start=4]
\item \textbf{Flash the microcontroller}
  \begin{numlist}
  \item Connect the microcontroller via USB.
  \item Run \code{usb_check}. This command makes sure that the WSL guest
    system can see the USB port.
  \item Start the upload with \code{make upload HEX=led_pin2.hex}
  \end{numlist}
\end{alphatasks}
\end{region}

\begin{region}{683.1}
\begin{alphatasks}[start=5]
\item \textbf{Measure with the oscilloscope}\\
  Use an oscilloscope and measure the signal at pin~2 of the
  microcontroller. You should see a square wave with 250\,ns HIGH and
  250\,ns LOW.
\end{alphatasks}
\end{region}

\nextpage
%% ===================================================================== 3 ==
\begin{region}{69.4}
\heading{Command overview for navigating the file system}
\end{region}

\pagedraw{
  \fill[tablegray] (56.9,93.7) rectangle (538.4,114.1);
  \fill[tablegray] (56.9,146.2) rectangle (538.4,166.6);
  \fill[tablegray] (56.9,187.0) rectangle (538.4,207.4);
  \fill[tablegray] (56.9,227.8) rectangle (538.4,271.8);
  \foreach \y in {73.3,93.7,114.1,146.2,166.6,187.0,207.4,227.8,271.8}
    \draw[tline] (56.4,\y) -- (538.9,\y);
  \foreach \x in {56.9,133.6,252.0,538.4}
    \draw[tline] (\x,72.8) -- (\x,272.3);
}

\begin{region}[x=61.0,width=68,free]{87.0}
\cell\bfseries Command
\end{region}
\begin{region}[x=256.1,width=278,free]{87.0}
\cell\bfseries Explanation
\end{region}

\begin{region}[x=61.0,width=68,free]{107.5}\cell pwd\end{region}
\begin{region}[x=137.8,width=110,free]{107.5}\cell print working directory\end{region}
\begin{region}[x=256.1,width=278,free]{107.5}\cell
Shows where you are in the file system
\end{region}

\begin{region}[x=61.0,width=68,free]{127.9}\cell ls\end{region}
\begin{region}[x=137.8,width=110,free]{127.9}\cell list directory content\end{region}
\begin{region}[x=256.1,width=278,free]{127.9}\cell
List of all files and subdirectories in the current directory
\end{region}

\begin{region}[x=61.0,width=68,free]{160.0}\cell cd\end{region}
\begin{region}[x=137.8,width=110,free]{160.0}\cell change directory\end{region}
\begin{region}[x=256.1,width=278,free]{160.0}\cell
Always changes into your own home directory
\end{region}

\begin{region}[x=61.0,width=68,free]{180.3}\cell cd ..\end{region}
\begin{region}[x=256.1,width=278,free]{180.3}\cell
Changes one directory up.
\end{region}

\begin{region}[x=61.0,width=68,free]{200.7}\cell cd foo\end{region}
\begin{region}[x=256.1,width=278,free]{200.7}\cell
Changes into the subdirectory foo (if it exists)
\end{region}

\begin{region}[x=61.0,width=68,free]{221.1}\cell cd foo/bar\end{region}
\begin{region}[x=256.1,width=278,free]{221.1}\cell
Changes into the sub-subdirectory foo/bar (if it exists)
\end{region}

\begin{region}[x=61.0,width=68,free]{241.5}\cell cat foo.txt\end{region}
\begin{region}[x=137.8,width=110,free]{241.5}\cell
concatenate files and print on the standard output
\end{region}
\begin{region}[x=256.1,width=278,free]{241.5}\cell
If the file foo.txt exists in the current directory, its content is
printed. Handy for quickly checking what is in a file.
\end{region}

\placefigure[187.7]{412.6}{301.9}{133.8}{flowchart-sbi-loop.tex}

\begin{region}{296.3}
\heading{Step 2: A second look at the assembly program}
\end{region}

\begin{region}[width=344,pitch=13.17]{323.9}
In this task, you will work out how the assembly program is structured --
and why it does the same thing as the machine code program that you wrote
by hand in the last lab session.
\end{region}

\begin{region}{389.8}
\begin{asmlisting}
\aN{1}         \aI{sbi}     \aK{0x0a}, \aK{2}     \aC{; 2 cycles}
\aN{2}
\aN{3 loop}:   \aI{sbi}     \aK{0x0b}, \aK{2}     \aC{; 2 cycles}
\aN{4}         \aI{nop}                 \aC{; 1 cycle}
\aN{5}         \aI{nop}                 \aC{; 1 cycle}
\aN{6}         \aI{cbi}     \aK{0x0b}, \aK{2}     \aC{; 2 cycles}
\aN{7}         \aI{rjmp}    \aR{loop}        \aC{; 2 cycles}
\end{asmlisting}
\end{region}

\begin{region}[width=344,pitch=12.1]{478.4}
Each line contains one machine instruction in symbolic form -- this is
assembly language. The instruction is on the far left, followed by its
arguments, if any. Everything after \code{;} is a comment; here it states
how many clock cycles the instruction takes.
\end{region}

\figcaption{444.3}{511.2}{Flowchart}

\begin{region}{539.6}
\subheading{The instructions in detail}
\end{region}

\pagedraw{
  \fill[tablegray] (56.9,564.3) rectangle (538.4,596.4);
  \fill[tablegray] (56.9,617.2) rectangle (538.4,637.9);
  \foreach \y in {543.5,564.3,596.4,617.2,637.9,668.2}
    \draw[tline] (56.4,\y) -- (538.9,\y);
  \foreach \x in {56.9,118.2,265.3,538.4}
    \draw[tline] (\x,543.0) -- (\x,668.7);
}

\begin{region}[x=61.0,width=54,free]{557.2}\cell\bfseries Instruction\end{region}
\begin{region}[x=122.3,width=139,free]{557.2}\cell\bfseries Meaning\end{region}
\begin{region}[x=269.5,width=265,free]{557.2}\cell\bfseries Effect in the program\end{region}

\begin{region}[x=61.0,width=54,free]{577.2}\cell\ttfamily sbi\end{region}
\begin{region}[x=122.3,width=139,free]{578.0}\cell Set Bit in I/O Register\end{region}
\begin{region}[x=269.5,width=265,free,pitch=12.1]{578.0}\cell
\makebox[78bp][l]{\ttfamily sbi 0x0a, 2}Set pin~2 as output\\
\makebox[78bp][l]{\ttfamily sbi 0x0b, 2}Set pin~2 HIGH
\end{region}

\begin{region}[x=61.0,width=54,free]{609.4}\cell\ttfamily cbi\end{region}
\begin{region}[x=122.3,width=139,free]{610.2}\cell Clear Bit in I/O Register\end{region}
\begin{region}[x=269.5,width=265,free]{610.2}\cell
\makebox[78bp][l]{\ttfamily cbi 0x0b, 2}Set pin~2 LOW
\end{region}

\begin{region}[x=61.0,width=54,free]{630.2}\cell\ttfamily nop\end{region}
\begin{region}[x=122.3,width=139,free]{630.9}\cell No Operation\end{region}
\begin{region}[x=269.5,width=265,free]{630.9}\cell
No effect, used to stretch the timing
\end{region}

\begin{region}[x=61.0,width=54,free,pitch=11]{650.9}\cell\ttfamily
rjmp\\loop
\end{region}
\begin{region}[x=122.3,width=139,free]{651.7}\cell Relative jump to a label\end{region}
\begin{region}[x=269.5,width=265,free]{651.7}\cell
Jumps back to the label {\ttfamily loop:}
\end{region}

\begin{region}{691.0}
Labels such as \code{loop:} serve as jump targets. The instruction
\code{rjmp loop} makes the program continue at line~3 -- this creates an
endless loop. The flowchart shows the sequence over time graphically. It can
help you understand how this program produces a square wave.
\end{region}

\begin{region}{752.0}
\subheading{Task}
Modify the assembly program so that a square wave with 500\,ns HIGH and
500\,ns LOW is generated at pin~2. As a reminder: one clock cycle takes
62.5\,ns.
\end{region}

\nextpage
%% ===================================================================== 4 ==
\begin{region}{69.4}
\heading{Step 3: Blinking more elegantly with exclusive OR}
\end{region}

\begin{region}[pitch=13.27]{97.0}
The following program -- as before -- makes sure that a symmetric square
wave is generated at pin~2: 250\,ns HIGH and 250\,ns LOW, alternating.\\
The difference lies in the code: this time we use registers and a new
instruction -- the exclusive OR (\code{eor}) in line~5.
\end{region}

\begin{region}{163.5}
\begin{asmlisting}
  \aN{1}         \aI{ldi}     \aR{r19},    \aK{4}           \aC{; 1 cycle}
  \aN{2}         \aI{out}     \aK{0x0a},   \aR{r19}         \aC{; 1 cycle}
  \aN{3}         \aI{ldi}     \aR{r18},    \aK{0}           \aC{; 1 cycle}
  \aN{4 loop}:
  \aN{5}         \aI{eor}     \aR{r18},    \aR{r19}         \aC{; 1 cycle}
  \aN{6}         \aI{out}     \aK{0x0b},   \aR{r18}         \aC{; 1 cycle}
  \aN{7}         \aI{rjmp}    \aR{loop}                \aC{; 2 cycles}
\end{asmlisting}
\end{region}

\begin{region}[pitch=14.3]{261.3}
\textbf{Before we analyze what exactly the code does, we first test whether
it really does what it promises.}
\end{region}

\begin{region}{304.1}
\subheading{Tasks}
\end{region}

\begin{region}[pitch=14.3]{318.4}
\begin{alphatasks}
\item In the editor \code{nano} (left terminal), save your current file with
  \textbf{Ctrl\,+\,O} (or \textbf{Ctrl\,+\,S}) and exit the editor with
  \textbf{Ctrl\,+\,X}.
\end{alphatasks}
\end{region}

\begin{region}{346.3}
\begin{alphatasks}[start=2]
\item Create a new file with \code{nano led_pin2_xor.s}
\end{alphatasks}
\end{region}

\begin{region}{359.5}
\begin{alphatasks}[start=3]
\item Type in the program shown above and save it.
\end{alphatasks}
\end{region}

\begin{region}{372.9}
\begin{alphatasks}[start=4]
\item Assemble and flash the new program with
  \code{make upload HEX=led_pin2_xor.hex}
\end{alphatasks}
\end{region}

\begin{region}{386.2}
\begin{alphatasks}[start=5]
\item Check with the oscilloscope whether pin~2 switches between HIGH and
  LOW as claimed.
\end{alphatasks}
\end{region}

\begin{region}{413.5}
\subheading{Explanation of the program (and one more task)}
\end{region}

\begin{region}[pitch=13.17]{427.1}
This program uses registers and an XOR operation (\code{eor}) to toggle the
state of pin~2. Registers are small storage cells in the processor. The
ATmega328P has 32 registers (r0 to r31), each 8 bits wide.
Unlike normal RAM, they can be used directly in arithmetic operations.
Besides addition and subtraction, an exclusive OR is also possible. The
following describes what the individual lines do.
\end{region}

\begin{region}[pitch=12.3]{507.1}
\textbf{Line 1:} \code{ldi r19, 4}\\
Writes the number 4 into register \code{r19}. Afterwards,
\code{r19 = b00000100}, i.e., only bit~2 is set.
\end{region}

\begin{region}[pitch=12.2]{535.4}
\textbf{Line 2:} \code{out 0x0a, r19}\\
Configures pins 0 to 7 of the microcontroller as inputs or outputs. Since
only bit~2 is set in \code{r19}, only pin~2 is enabled as an output.
\end{region}

\begin{region}[pitch=12.4]{575.7}
\textbf{Line 3:} \code{ldi r18, 0}\\
Initializes \code{r18} with 0, i.e., all bits are cleared.
\end{region}

\begin{region}[pitch=12.2]{604.1}
\textbf{Line 5:} \code{eor r18, r19}\\
Combines \code{r18} with \code{r19} using an exclusive OR. This switches
bit~2 in \code{r18} -- this is also called ``toggling'': if it was 0
before, it becomes 1; if it was 1 before, it becomes 0.\\
The content of \code{r18} therefore alternates between \code{b00000000} and
\code{b00000100}.
\end{region}

\begin{region}[pitch=12.2]{656.7}
\textbf{Line 6:} \code{out 0x0b, r18}\\
Outputs the content of \code{r18} to pins 0 to 7 of the microcontroller.\\
If a bit in \code{r18} is set and the corresponding pin is configured as an
output, this pin is set HIGH. If the bit is not set, the pin (provided it is
configured as an output) is set LOW.
\end{region}

\begin{region}{730.3}
\textbf{Task:} Explain to us why the exclusive OR in line~5 ``toggles''
bit~2.
\end{region}

\nextpage
%% ===================================================================== 5 ==
\begin{region}{69.4}
\heading{Step 4: Simple delay with one register}
\end{region}

\begin{region}[pitch=13.53]{95.9}
Compared with the previous program, only lines 7 to 10 are new in this
example. So you do \textbf{not have to type in the entire program again}.
As you will measure in a moment, these additional lines make pin~2 switch
between HIGH and LOW more slowly.
\end{region}

\begin{region}{153.0}
\begin{asmlisting}
  \aN{1}         \aI{ldi}     \aR{r19},    \aK{4}      \aC{; 1 cycle}
  \aN{2}         \aI{out}     \aK{0x0a},   \aR{r19}    \aC{; 1 cycle}
  \aN{3}         \aI{ldi}     \aR{r18},    \aK{0}      \aC{; 1 cycle}
  \aN{4 loop}:
  \aN{5}         \aI{eor}     \aR{r18},    \aR{r19}    \aC{; 1 cycle}
  \aN{6}         \aI{out}     \aK{0x0b},   \aR{r18}    \aC{; 1 cycle}
  \aN{7}         \aI{ldi}     \aR{r24},    \aK{0x02}   \aC{; 1 cycle (set counter for countdown)}
  \aN{8 delay}:
  \aN{9}         \aI{subi}    \aR{r24},    \aK{1}      \aC{; 1 cycle}
 \aN{10}         \aI{brne}    \aR{delay}          \aC{; 2 cycles if taken, 1 otherwise}
 \aN{11}
 \aN{12}         \aI{rjmp}    \aR{loop}           \aC{; 2 cycles}
\end{asmlisting}
\end{region}

\begin{region}{304.7}
\subheading{Tasks}
\end{region}

\begin{region}[pitch=13.5]{318.2}
\begin{numlist}
\item Exit the editor. Then run the command
  \code{cp led_pin2_xor.s led_pin2_delay.s}. This gives you a copy of the
  file \code{led_pin2_xor.s} named \code{led_pin2_delay.s}.
\end{numlist}
\end{region}

\begin{region}{345.3}
\begin{numlist}[start=2]
\item Add lines 7 to 10 from the listing above to \code{led_pin2_delay.s}.
\end{numlist}
\end{region}

\begin{region}[pitch=13.3]{358.5}
\begin{numlist}[start=3]
\item The figure on the next page shows a flowchart for the program. You can
  see that there are now two nested loops. The inner loop is new:
  \begin{steps}
  \item In line~7, register \code{r24} is initialized with the value
    \code{0x02} (i.e., \code{2}).
  \item Lines 8 to 10 form the inner loop, in which \code{r24} is counted
    down until it has reached the value \code{0}. Only then is the loop
    exited.
  \end{steps}
  The project companion derives the following:
\end{numlist}
\end{region}

\begin{region}[x=74.6,right=0]{452.7}
\hltext{If {\ttfamily r24} is initialized with \textbf{n}, one pass of the outer
loop takes exactly \textbf{3·n + 4} clock cycles.}
\end{region}

\begin{region}[x=74.6,right=0,pitch=13.3]{480.1}
\hltext{This means: pin~2 stays HIGH for \textbf{3·n + 4} clock cycles and then
LOW for just as many cycles before the state switches again.}
\end{region}

\begin{region}[x=74.6,right=0,pitch=14.1]{519.8}
You will now verify this claim experimentally. To do so, fill in the table
below. For different values of \textbf{n}, you
\end{region}

\begin{region}[x=74.7,right=0,pitch=13.25]{548.4}
\begin{dotlist}
\item adjust the initialization of \code{r24} in line~7,
\item flash the program to the microcontroller again,
\item and then check with the oscilloscope how long pin~2 is HIGH and LOW.
\end{dotlist}
\end{region}

\begin{region}{596.0}
\subheading{Table for the measurements (1 clock cycle takes 62.5\,ns)}
\end{region}

\pagedraw{
  \foreach \y in {633.2,674.0,714.8}
    \fill[tablegray] (56.9,\y) rectangle (538.4,\y+20.4);
  \foreach \y in {601.1,633.2,653.6,674.0,694.4,714.8,735.2,755.6}
    \draw[tline] (56.4,\y) -- (538.9,\y);
  \foreach \x in {56.9,84.4,151.2,311.9,439.3,538.4}
    \draw[tline] (\x,600.6) -- (\x,756.1);
}

\begin{region}[x=56.9,width=27.5,free]{614.9}\cell\centering n\end{region}
\begin{region}[x=88.4,width=58.8,free]{614.7}\raggedleft 3·n + 4\end{region}
\begin{region}[x=156.3,width=152,free]{614.9}\cell
Exact duration for 3·n + 4 cycles
\end{region}
\begin{region}[x=311.9,width=127.4,free]{614.9}\cell\centering
Measured\\HIGH duration
\end{region}
\begin{region}[x=439.3,width=99.1,free]{614.9}\cell\centering
Measured\\LOW duration
\end{region}
\begin{region}[x=56.9,width=27.5,free]{647.0}\cell\centering 1\end{region}
\begin{region}[x=56.9,width=27.5,free]{667.4}\cell\centering 2\end{region}
\begin{region}[x=56.9,width=27.5,free]{687.8}\cell\centering 3\end{region}
\begin{region}[x=56.9,width=27.5,free]{708.2}\cell\centering 4\end{region}
\begin{region}[x=56.9,width=27.5,free]{728.6}\cell\centering 255\end{region}
\begin{region}[x=56.9,width=27.5,free]{749.0}\cell\centering 0\end{region}

\nextpage
%% ===================================================================== 6 ==
\placefigure[539.7]{123.4}{151.1}{348.5}{flowchart-delay.tex}

\nextpage
%% ===================================================================== 7 ==
\begin{region}{69.4}
\heading{Step 5: 32-bit delay with four registers}
\end{region}

\begin{region}[pitch=13.2]{96.2}
A single register of our microprocessor can only store values between
\code{0x00} (i.e., 0) and \code{0xFF} (i.e., 255). That is not enough to
delay the outer loop so much that the blinking of an LED connected to
pin~2 can be perceived with the naked eye.
\end{region}

\begin{region}[pitch=13.4]{149.2}
The following program therefore uses a 32-bit counter that is implemented
with the four registers \code{r24} to \code{r27}. Together, these four
registers form a so-called register group, which is treated like a single
large register. The arrangement is:
\end{region}

\begin{region}[x=59.8,right=0]{202.6}
\code{r27 : r26 : r25 : r24}
\end{region}

\begin{region}[pitch=13.4]{229.0}
This group stores a 32-bit value in the so-called little-endian format --
that is, the least significant byte is in \code{r24}, the most significant
one in \code{r27}.
\end{region}

\begin{region}{268.8}
Example: if the registers are set as follows
\end{region}

\begin{region}{295.4}
\code{r24 = 0x12, r25 = 0x34, r26 = 0x56, r27 = 0x78}
\end{region}

\begin{region}{321.8}
then the combined value is:
\end{region}

\begin{region}{348.4}
\code{n = r27:r26:r25:r24 = 0x78563412}
\end{region}

\begin{region}{374.9}
In the program below, the registers are initialized as follows:
\end{region}

\begin{region}{401.5}
\code{r24 = 0x00, r25 = 0x02, r26 = 0x00, r27 = 0x00}
\end{region}

\begin{region}{427.9}
This corresponds to:
\end{region}

\begin{region}{454.5}
\code{n = 0x00000200 = 512}
\end{region}

\begin{region}{480.9}
That is: the inner loop runs 512 times.
\end{region}

\begin{region}{507.5}
\begin{asmlisting}
  \aN{1}         \aI{ldi}     \aR{r19},    \aK{4}         \aC{; for pin 2 (4 = 2\textasciicircum2)}
  \aN{2}         \aI{out}     \aK{0x0a},   \aR{r19}
  \aN{3}         \aI{ldi}     \aR{r18},    \aK{0}
  \aN{4 loop}:
  \aN{5}         \aI{eor}     \aR{r18},    \aR{r19}       \aC{; 1 cycle}
  \aN{6}         \aI{out}     \aK{0x0b},   \aR{r18}       \aC{; 1 cycle}
  \aN{7}
  \aN{8}         \aI{ldi}     \aR{r24},    \aK{0x00}      \aC{; 1 cycle}
  \aN{9}         \aI{ldi}     \aR{r25},    \aK{0x02}      \aC{; 1 cycle}
 \aN{10}         \aI{ldi}     \aR{r26},    \aK{0x00}      \aC{; 1 cycle}
 \aN{11}         \aI{ldi}     \aR{r27},    \aK{0x00}      \aC{; 1 cycle}
 \aN{12 delay}:
 \aN{13}         \aI{subi}    \aR{r24},    \aK{1}         \aC{; 1 cycle}
 \aN{14}         \aI{sbci}    \aR{r25},    \aK{0}         \aC{; 1 cycle}
 \aN{15}         \aI{sbci}    \aR{r26},    \aK{0}         \aC{; 1 cycle}
 \aN{16}         \aI{sbci}    \aR{r27},    \aK{0}         \aC{; 1 cycle}
 \aN{17}         \aI{brne}    \aR{delay}             \aC{; 2 cycles if taken, 1 otherwise}
 \aN{18}
 \aN{19}         \aI{rjmp}    \aR{loop}              \aC{; 2 cycles}
\end{asmlisting}
\end{region}

\nextpage
%% ===================================================================== 8 ==
\begin{region}{69.4}
\subheading{Tasks}
\end{region}

\begin{region}[pitch=13.37]{95.9}
\begin{numlist}
\item You do not have to type in the entire program again.\\
  Compared with \code{led_pin2_delay.s}, only lines 9 to 11 (initialization
  of the registers) and lines 14 to 16 (decrementing the register group)
  are new.\\
  {\aslabarrow} Adjust your file \code{led_pin2_delay.s} accordingly.
\end{numlist}
\end{region}

\begin{region}[pitch=13.2]{149.3}
\begin{numlist}[start=2]
\item The project companion explains in detail how to decrement a register
  group correctly. Here you are to verify the following claim
  experimentally (it only holds for n > 0):
\end{numlist}
\end{region}

\begin{region}[x=74.7,right=0]{189.7}
\hltext{If {\ttfamily r24} is initialized with \textbf{n}, one pass of the outer
loop takes exactly \textbf{6·n + 7} clock cycles.}
\end{region}

\begin{region}[x=74.7,right=0]{217.2}
\hltext{It follows that pin~2 stays HIGH for \textbf{6·n + 7} clock cycles and
then LOW for just as long.}
\end{region}

\begin{region}[x=74.7,right=0]{243.7}
Fill in this table (1 clock cycle takes 62.5\,ns):
\end{region}

\begin{region}{284.1}
\begin{numlist}[start=3]
\item \textbf{Now calculate the value n at which pin~2 blinks at a rate of
  exactly 1\,Hz.}
\end{numlist}
\end{region}

\pagedraw{
  \fill[tablegray] (87.2,325.1) rectangle (529.3,345.5);
  \foreach \y in {293.1,325.1,345.5,365.9}
    \draw[tline] (86.7,\y) -- (529.8,\y);
  \foreach \x in {87.2,112.5,173.8,332.8,438.4,529.3}
    \draw[tline] (\x,292.6) -- (\x,366.4);
}

\begin{region}[x=87.2,width=25.3,free]{306.8}\cell\centering n\end{region}
\begin{region}[x=116.5,width=53.3,free]{306.6}\raggedleft 6·n + 7\end{region}
\begin{region}[x=178.1,width=152,free]{306.8}\cell
Exact duration for 6·n + 7 cycles
\end{region}
\begin{region}[x=332.8,width=105.6,free]{306.8}\cell\centering
Measured\\HIGH duration
\end{region}
\begin{region}[x=438.4,width=90.9,free]{306.8}\cell\centering
Measured\\LOW duration
\end{region}
\begin{region}[x=87.2,width=25.3,free]{338.9}\cell\centering 1\end{region}
\begin{region}[x=87.2,width=25.3,free]{359.3}\cell\centering 2\end{region}

\begin{region}[x=74.7,right=0,pitch=13.1]{389.7}
For this to work, the pin must be HIGH for 500\,ms and LOW for 500\,ms. To
do so, solve the equation:
\end{region}

\begin{region}[x=74.7,right=0]{429.2}
(6·n + 7) × 62.5\,ns = 500\,ms = 500000000\,ns
\end{region}

\begin{region}[x=74.7,right=0]{456.2}
\begin{arrowlist}
\item Calculate the matching value for n.
\end{arrowlist}
\end{region}

\begin{region}[x=74.7,right=0]{469.4}
\begin{arrowlist}
\item If n is not an integer, round it to an integer.
\end{arrowlist}
\end{region}

\begin{region}[x=74.7,right=0,pitch=12.5]{482.5}
\begin{arrowlist}
\item Initialize the counter accordingly, flash the program, and check with
  the oscilloscope whether the signal matches your expectation.
\end{arrowlist}
\end{region}

\begin{region}[x=74.7,right=0,pitch=12.5]{508.9}
\begin{arrowlist}
\item Additionally, connect an LED with a series resistor to pin~2 and
  watch it blink at 1\,Hz -- hopefully pleasantly slowly and evenly.
\end{arrowlist}
\end{region}

\end{document}
